\author{OTS}
\documentclass[fontsize=16pt]{mybib}
\renewcommand{\DFA}{12mm}
\renewcommand{\DZA}{1.3}

\graphicspath{{../../../assets/images/}}
\setheaderlogo{mnr-green.png}

% Absatzformatierung
\setlength{\parindent}{0pt}
\setlength{\parskip}{0.6em}
\renewcommand{\arraystretch}{0.8}

\begin{document}
    \begin{enumerate}
        \item Das Problem \biblezit{Neh}{1:1-4}
        \item Hilfe erbeten \biblezit{Neh}{1:5-11}
        \item Die Erhörung \biblezit{Neh}{2:1-8}
    \end{enumerate}

    \begin{enumerate}
        \item Das Problem
        \begin{enumerate}
            \item Vers 1: Zeit und Vorwort
            \item Vers 2: Die Informanten
            \item Vers 3: Die Information
            \item Vers 4a: Resultat der Information
        \end{enumerate}
        \item Für Hilfe beten
        \begin{enumerate}
            \item Vers 4b: Neh wird tief betrübt
            \item Vers 5: Gott loben
            \item Vers 6: Neh erniedrigt sich vor Gott
            \item Vers 7: Neh bekennt seien Fehler
            \item Vers 8-9: Neh erinner Gott an seine Zusagen
            \item Vers 10: Sie gehören dazu
            \item Vers 11: Die Bitte an Gott
        \end{enumerate}
        \item Die Gebetserhörung
        \begin{enumerate}
            \item Vers 1: Dauer bis zur Erhörung
            \item Vers 2: Der König spricht
            \item Vers 3: Neh nennt den Grund
            \item Vers 4: Die nachfrage des Königs
            \item Vers 5: Neh sagt sein anliegen
            \item Vers 6: Der König will Details wissen
            \item Vers 7-8: Die Details
        \end{enumerate}
    \end{enumerate}

\end{document}